\begin{table}[!htbp]
\centering
\caption{Área relativa dos footprints de bom desempenho nos quatro modelos.}
\label{tab:res-footprints}
\begin{singlespace}
\normalsize
\setlength{\tabcolsep}{7pt}
\renewcommand{\arraystretch}{1.12}
\begin{tabular}{lrrrr}
\hline\hline
\noalign{\vskip 3pt}
\textbf{Método} & \textbf{\shortstack{55\\controle}} & \textbf{\shortstack{55\\topologia}} & \textbf{\shortstack{74\\controle}} & \textbf{\shortstack{74\\topologia}} \\
\hline
IBS & 100,00 & 100,00 & 100,00 & 100,00 \\
MIP & 15,64 & 39,82 & 54,47 & 26,33 \\
CUTH & 0,00 & 0,00 & 4,38 & 11,07 \\
LBBD & 15,64 & 39,82 & 38,78 & 28,93 \\
ILS & 15,64 & 39,82 & 38,78 & 28,93 \\
RS & 0,00 & 0,00 & 4,38 & 6,47 \\
\hline\hline
\end{tabular}
\par\smallskip
{\small Nota: área do footprint dividida pela área de referência construída pelo TRACE no próprio mapa, em porcentagem. As regiões usam classes previstas, não rótulos observados. Zero indica ausência de região construída, não prova de incapacidade do método. As áreas absolutas de mapas diferentes não são diretamente comparáveis.\par}
{\fonte{elaboração própria a partir dos resultados dos quatro modelos.}}
\end{singlespace}
\end{table}
